\chapter{}

Pietro tornava solo da lunghissime passeggiate in campagna, dopo essersi consigliato anche con l'aria. Talvolta gli era parso impossibile che Ghìsola avesse amato qualcuno, perchè sarebbe stato una contaminazione della sua bellezza. Piuttosto era lui un geloso!

Talvolta si diceva: «Sono proprio a Siena? Non mi pare la stessa. Certamente, il suo cielo ora è più azzurro di prima: non era così una volta». Notò che d'estate, verso sera, nella Piazza del Campo rimane una luce pallida e tepida, un avanzo del meriggio; simile alla luce d'una lanterna, che illumini soltanto là dentro; mentre le persone, che attraversano quello spazio, sembrano lontane nel tempo, con un silenzio indefinibile.

«Quando ci sarà anche Ghìsola, le dirò quel che provo.»

Tutte le mattine si svegliava con un sospiro. E come si ricordava bene dei sogni!

Ma senza Ghìsola non poteva vivere; e, verso la metà d'agosto, decise d'andare a prenderla, perchè tornasse a Radda ad aspettare il loro matrimonio; un anno forse, un anno e mezzo al più. Perchè non avrebbe avuto il consenso? Intanto, facendola stare a Radda, si sentiva più sicuro di lei.

Da Rebecca si fece prestare il denaro per il viaggio.

Ma a Firenze, in quelle poche ore, gli pareva d'essere sempre a Siena, in cima alla via di Camporegio, dove era andato tutti i giorni quando faceva la scuola tecnica. È breve la distanza tra la mole rude e rossiccia di San Domenico e le case che s'arrampicano alla rinfusa, un'altra volta, in ogni direzione attorno al Duomo, fermandovisi sotto a pena che lo toccano; ma, a guardare di lì la profondità vuota di Fontebranda, ci si sente mozzare il respiro.

L'Ospedale, alto su le mura, rosso sangue, lo vedeva doventare del colore della terra bruciata; il turchino del cielo, bigio. E poi le prime stelle, qua e là, così sparse che gli facevano angoscia.

I vicoli, simili a spaccature e a cretti enormi, s'anneravano.

Tra i giardini e gli orti, l'uno più alto dell'altro, chiusi dentro i muri rettangolari, che spesso hanno a comune, nelle insenature o nelle sporgenze delle colline, e seguendo i loro pendii diseguali, il barlume della notte gli sembrava che cadesse come quando piove a dirotto.

Un briaco cominciava a cantare e poi smetteva. La Costaccia come il parapetto d'un abisso, e il Costone quasi a picco, con il suo arco greve e largo che lo tiene fermo perchè sopra ci passi un'altra strada, salgono di squincio verso le case.

Non due tetti della stessa altezza, anche se accanto. Grumoli piccoli e grandi di case che s'allungano parallelamente obliqui e storti: alcune volte le case stanno a due e tre angoli l'uno dentro l'altro, a cerchio, a nodi, serrate insieme, mescolate, aggrovigliate, con curve rotte o schiacciate, sempre con improvvisi cambiamenti; obbedendo alle forme delle colline, ai pendii e alle svolte delle vie, alle piazze che dall'alto paiono buche.

Ad un tratto, uno stacco tra due case, e poi le altre che s'afferrano e si tengono

ancora, con forza, pigiandosi e abbassandosi e poi risalendo e girando per sparire leste leste dietro quelle che hanno un movimento affatto diseguale e che vengono incontro dalla parte opposta; salite su; ma anche queste s'interrompono quasi subito per doventare una raggiera più larga, irregolare, tutta piana oppure contorta; dentro la quale si mettono e s'avventano case, di sghembo, a traverso, come riescono e possono; spinte da altre che fanno l'effetto di volersi accomodare meglio ed assestarsi, ciascuna per conto proprio.

Le case, bassissime, quasi per affondare nella campagna, da Porta Ovile, da Fontebranda, dai Tufi, sorreggono quelle che hanno a ridosso, le trattengono dalla loro voglia di sparpagliarsi più rade; i punti più alti sono come richiami alle case costrette ad obbedire per non restare troppo sole.

Nei rialzi sembra che ci sia un parapiglia a mulinello, negli abbassamenti le case precipitano l'una addosso all'altra; come frane. Oppure si possono contare fino a dieci file di tetti, lunghe lunghe, sempre più alte: di fianco, altre file che vanno in senso perpendicolare alle prime.

La Torre del Mangia esce fuori placida da tutto quell'arruffio.

E attorno alla città, gli olivi e i cipressi si fanno posto tra le case; come se, venuti dalla campagna, non volessero più tornare a dietro.

Ma gli pareva d'essere inseguito da suo padre, pur sentendosi rasserenato dal campanile di Giotto, da Santa Maria del Fiore, da quelle strade che conosceva, già percorse in quella specie di perdizione sempre più accanita. Aveva voglia di riparlare con qualcuno dei suoi compagni, di spiegare a loro l'equivoco avuto, e come si fosse perso per una ragione che non sapeva dire; per quanto gli dispiacesse tenere segreti anche ora che sentiva la necessità squisita d'aver qualche cosa da nascondere; una cosa che forse era come la sua anima stessa.

Un venditore di limoni, sotto un ombrello verde con le stecche di legno, era seduto al principio del Ponte alle Grazie. Qualche facchino e qualche persona indefinibile sonnecchiavano appoggiati al muricciolo dell'argine.

Un'allodola volò dagli alberi di San Miniato, verso le Cascine, come una cosa scintillante.

Andando verso la Piazza della Signoria, fresca e annaffiata, si cominciava a rivedere la gente: più fitta in Via Calzaioli e nella Piazza del Duomo. In fondo a Via Cavour, il poggio di Fiesole; alto e verde. A Badia, quando scese dal tranvai, Pietro arrossì quantunque non ci fosse nessuno. E scrutò sotto le persiane, per scorgervi qualche viso che guardasse nella strada: soltanto piante di geranii polverosi.

Apertogli l'uscio proprio da Ghìsola, che però non lo fece entrare, egli subito si dolse che non fosse già andata a Radda; ed ella rispose che aspettava lui e voleva prima esser sicura che i suoi genitori l'avebbero volentieri ripresa in casa.

Gli era inesplicabile la sensazione di trovarsi con lei già da tanto tempo.

— E perchè no? Sono cattivi con te?

— Io non ci sto volentieri.

Gli fece caso che rispondesse proprio a quel modo e non altrimenti. L'accarezzò, pregandola:

— Tu non mi devi rispondere di no; devi aspettare a casa tua. Mi farai piacere.

Poi pensò: «Perchè le domando di fare così?».

— Se tu vuoi....

Visto ch'ella era per ubbidire, chiese:

— Vieni a Siena con me, allora.

Ella sorrise e gli fece cenno di tacere. Era convinto che dovesse provare una gran dolcezza ad ubbidirgli; ma Ghìsola, che aveva voglia di scherzare più che d'altro, gli chiese: — Ti piaccio meno?

— Perchè dovresti piacermi meno?

E le accarezzò tutta la faccia: ella si discosto e gli guardò la punta delle dita.

— Perchè non vuoi? Ti aspetto nella strada, verso la Badia.

— Verrò. Ora vattene.

Le baciò ambedue le mani, tenendogliele insieme, mentre ella si tirava a dietro, quasi chiudendogli l'uscio in faccia.

Ed egli pensava, scendendo le scale: «Ha sofferto. Soffre perchè deve stare in una casa che non è sua. I genitori, forse, non le hanno più scritto; i parenti la invidiano. M’è parsa più sensuale; ma io devo rispettarla lo stesso, anzi di più: dopo, la odierei.»

Invece non gli fece caso che potesse venirsene via così a pena glie ne aveva parlato.

Il signor Alberto s’era impigliato in un processo di fallimento; e da una quindicina di giorni non si faceva più vedere da nessuno, nè meno da lei, che andava a trovarlo, di rado, qualche mezz’ora, nello studio d’uno dei suoi avvocati dove ormai passava tutto il suo tempo. Egli l’aveva pregata di tornare a Radda, soltanto finchè il processo non fosse finito; anche perchè i parenti della moglie, ch’erano tra i testimonii, non soffiassero nella brace.

Denari non li dava più; e, più d’una volta, Ghìsola aveva dovuto cominciare a contentarsi di pane mangiato soltanto con qualche frutta. Ma, non volendo tornare a casa e non avendo dove andare, aspettava prima di decidere qualche cosa.

Così non aveva, dunque, dopo l’arrivo di Pietro, che da incaricare Beatrice di salutare il suo amico, pregandolo che non la dimenticasse.

Tuttavia, per farle ricordare che Pietro l’aspettava, ci vollero le altre persuasioni di Beatrice; alla quale, evidentemente, il padrone aveva ricorso anche per questa faccenda.

La donna l’abbracciò piangendo; con una tenerezza che la fece sorridere, lacrimando.
